\begin{proof}
\emph{(i)} Put $\chi = T^{-f/2}Z$. Then $\ddot\chi + (1 + f)\dot\chi/T = T^{-f/2}\bigl(\ddot Z + \dot Z/T - \tfrac14 f^{2}Z/T^{2}\bigr)$, so (E) reads $\ddot Z + \dot Z/T + (k^{2}T^{-2P} - \tfrac14 D\,T^{-2})Z = 0$. With $s = kT^{q}/|q|$ one has $T\,d/dT = q\,s\,d/ds$ and $k^{2}T^{-2P} = q^{2}s^{2}T^{-2}$, so $T^{2}(\ddot Z + \dot Z/T) = q^{2}(s^{2}Z_{ss} + sZ_{s})$ and the equation is Bessel's, $s^{2}Z_{ss} + sZ_{s} + (s^{2} - \nu^{2})Z = 0$ with $\nu^{2} = D/(4q^{2})$; this is the reduction of Fact \claimref{synthesis.imported-bessel-reduction-and-small-argument-hankel-forms} with $a = -f/2$ and $c = q$ after the shift of $\kappa$ into the order. The Frobenius form of the solutions, $T^{-f/2}\sum_{m}c_{m}s^{\pm\nu + 2m}$ with $c_{m} = -c_{m-1}/(4m(m \pm \nu))$, is verified term by term in exact arithmetic by the first check, which also confirms that a different order fails at the lowest term. By Theorem \claimref{synthesis.kasner-approach-coupled-modes-reach-measured-tilt-only-by-a-fitted-curvature-coupling}(iii) the exited spectrum of a mode in its adiabatic vacuum has index $3 - 2\nu$ for $\nu > 0$: the normalisation is free of $k$ and the exited branch carries $s^{-\nu}$, hence $k^{-\nu}$. The same holds beyond $P < 1$: the Wronskian of (E) is conserved up to the factor $T^{-(1 + f)}$, and for $\chi = N\,T^{-f/2}$ times a Hankel function of $s$ it equals $|N|^{2}\,q\,T^{-(1 + f)}$ times a constant free of $k$, since $T\,ds/dT = qs$ and the Wronskian of the Hankel pair in $s$ is a constant over $s$; so $N$ is free of $k$ whenever the kinetic weight of the mode's action is $T^{1 + f}$ with a coefficient free of $k$, for either sign of $q$ and for the de Sitter values of (v). So $n - 1 < 0$ exactly when $4\nu^{2} > 9$, that is $D > 9q^{2}$, with equality for scale invariance. With the measured value of Fact \claimref{synthesis.imported-planck-2018-lcdm-parameters}, the window $N_{\mathrm{obs}} \pm 3\sigma$ is $2\nu \in [3.0225, 3.0477]$, that is $3.0225^{2}q^{2} \le D \le 3.0477^{2}q^{2}$.

\emph{(ii)} On the derived metric $dT^{2} - T^{2p}\sum_{j \le d}(dx^{j})^{2}$ the volume factor is $T^{dp}$ and the equation of a massless minimally coupled field is $\ddot\chi + dp\,\dot\chi/T + k^{2}T^{-2p}\chi = 0$: $f = dp - 1$, $\kappa = 0$, $P = p$. Then $D - 9q^{2} = (dp - 1)^{2} - 9(1 - p)^{2}$ is a difference of squares, $(dp - 1 - 3 + 3p)(dp - 1 + 3 - 3p)$, which is the stated product; at $d = 3$ the second factor is $2$ and the first $6p - 4$. The order $|3p - 1|/(2|1 - p|)$ equals $|\tfrac32 + 1/(p - 1)|$, the form of Fact \claimref{synthesis.imported-wands-duality-of-perturbation-spectra}, whose statement of the index covers $p > 1$ as well. The Hubble rate is $p/T$ and the physical momentum $kT^{-p}$; their ratio $pT^{p - 1}/k$ grows toward $T = 0$ for $p < 1$ and toward large $T$ for $p > 1$. For $\tfrac13 \le p < 1$, $n - 1 = 3 - (3p - 1)/(1 - p) = (4 - 6p)/(1 - p)$, so $p = (4 - N)/(6 - N)$ at index $N$; for $p > 1$, $n - 1 = 3 - (3p - 1)/(p - 1) = -2/(p - 1)$, so $p = 1 - 2/N$. Evaluating at $N = -0.0351$ and at $N = -0.0351 \pm 0.0126$ gives the centres and intervals stated (rounded outward). With $p = 2/(3(1 + w_{\mathrm{f}}))$, $4 - 6p = 4w_{\mathrm{f}}/(1 + w_{\mathrm{f}})$ and $1 - p = (1 + 3w_{\mathrm{f}})/(3(1 + w_{\mathrm{f}}))$, so $n - 1 = 12w_{\mathrm{f}}/(1 + 3w_{\mathrm{f}})$.

\emph{(iii)} With exponents summing to one the volume factor is $T$, so the friction term is $\dot\chi/T$ and $f = 0$ (Theorem \claimref{synthesis.kasner-approach-gives-horizon-exit-and-index-three}(ii)); with $\kappa = 0$, $D = 0$. Equal exponents over $d$ directions with unit trace are $p = 1/d$, and $1/d < \tfrac23$ for $d \ge 2$. Unit trace is part of the setting of the cited theorems; for vacuum walls it is Theorem \claimref{dark-energy.frozen-point-datum-kasner-wall}, for scalar walls Theorem \claimref{synthesis.scalar-kasner-wall-is-reduced-five-dimensional-vacuum-wall}, and for stiff matter Theorem \claimref{synthesis.kasner-wall-matter-maxwell-bounces-stiff-moves-residue}. The statement is about these walls, not an audit of every field in the store.

\emph{(iv)} The rows are the stored ones: the cross modes have $f = 2(P_{b} - P_{a})$ (Theorem \claimref{synthesis.kasner-approach-coupled-modes-reach-measured-tilt-only-by-a-fitted-curvature-coupling}(i), (ii)), zero for an observed pair on the segment; a scalar with weight $b^{\lambda}$ has friction coefficient $1 + (\lambda - 1)P_{h}$ (Theorem \claimref{synthesis.kasner-hidden-scale-coupled-scalars-stay-blue}(i)), so $f = (\lambda - 1)P_{h}$, and for $\lambda \in \{-1, 0, 1\}$, $f^{2} \le 4P_{h}^{2}$. With $P_{h} = 1 - 3p$ and $q = 1 - p$, $9q^{2} - 4P_{h}^{2} = (3 - 3p - 2 + 6p)(3 - 3p + 2 - 6p) = (1 + 3p)(5 - 9p)$, positive for $p < \tfrac59$, hence on $[0, \tfrac12]$. For the Kaluza-Klein vector $f = 2(P_{h} - p) = 2(1 - 4p)$ and $f^{2} - 9q^{2} = (2 - 8p - 3 + 3p)(2 - 8p + 3 - 3p) = (1 + 5p)(11p - 5)$, positive exactly for $p > \tfrac{5}{11}$; and $16q^{2} - f^{2} = (4 - 4p - 2 + 8p)(4 - 4p + 2 - 8p) = 12(1 + 2p)(1 - 2p) \ge 0$ on the segment, so $\nu \le 2$ and the field strength's index $5 - 2\nu$ (clause (iii) of the same theorem on coupled modes) is at least $1$. Which points of the segment stored sources reach is clause (v) of that theorem. At the isotropic point $4\nu^{2} = (\lambda - 1)^{2} + 6\xi$ is the formula of Theorem \claimref{synthesis.tree-level-curvature-couplings-stay-blue-at-the-isotropic-kasner-point}(iv); over its twelve couplings (largest $\tfrac12$) and $|\lambda - 1| \le 2$ the largest value is $4 + 3 = 7$. At $\lambda = 1$ the window of (i) is $6\xi \in [3.0225^{2}, 3.0477^{2}]$, which gives the stated decimals exactly.

\emph{(v)} With scale factor $1/(HT)$, $T$ minus the conformal time, the Klein-Gordon equation of a field of mass $m$ in the unrescaled field variable is $\chi_{TT} - 2\chi_{T}/T + (k^{2} + m^{2}/(H^{2}T^{2}))\chi = 0$: the friction is twice the logarithmic derivative of the scale factor, $-2/T$, and the mass term is $m^{2}$ times the squared scale factor. So $1 + f = -2$, $P = 0$, $\kappa = -m^{2}/H^{2}$, and $D = 9 - 4m^{2}/H^{2} = 9 + 4w(w + 3) = (2w + 3)^{2}$ by the dictionary of Corollary \claimref{mass.weight-polynomial-two-series}. For $w \ge -\tfrac32$, $3 - 2\nu = -2w$. On $[-3, 0]$, $|2w + 3| \le 3$ with equality at the ends only. $D > 9$ is $w(w + 3) > 0$, that is $m^{2} < 0$. The index $-0.0351$ needs $-2w = -0.0351$ on the branch $w > 0$, $w = \tfrac{351}{20000}$, and then $m^{2}/H^{2} = -w(w + 3)$ is the stated fraction, the number of Proposition \claimref{synthesis.frozen-phase-spectrum-not-red}(iii). That the observed modes are inside the frozen point's Hubble length is Theorem \claimref{synthesis.frozen-point-geometry-reads-descent-as-vacuum}(ii).

\emph{(vi)} Orders $\tfrac12$ and $\tfrac32$ are (ii) at $p = \tfrac12$ and $\tfrac23$; the absence of exit is (ii) for an expansion with $p < 1$ and, for the stored descended phase, Proposition \claimref{synthesis.finite-range-descent-deposit-cannot-seed-cmb}(ii). Under $(a, k, A) \to (ca, ck, c^{3}A)$ at fixed $\Lambda$ and fixed time, $Aa^{-3}$ and $k/a$ are unchanged, so $\dot a/a$ and the mode equation are unchanged, and $ca$ solves the background equation with $c^{3}A$. A monomial $A^{r}\Lambda^{t}k^{u}$ has scaling weight $3r + u$ and, since $A$ and $\Lambda$ are inverse squared times and $k$ (with $a$ dimensionless) an inverse time, time dimension $-(2r + 2t + u)$; the two conditions have rank two, and the solutions are the multiples of $(2, 1, -6)$. On the dust stage a mode leaves the Hubble length when $k^{2} = a^{2}Aa^{-3}$, $a = A/k^{2}$, where the ratio of $\Lambda/3$ to $Aa^{-3}$ is $\Lambda A^{2}/(3k^{6}) = x/3$. A dimensionless spectrum built from $A$, $\Lambda$, $k$ alone and invariant under the scaling is therefore a function of $x$; at $\Lambda = 0$ it is the dust value, which is free of $k$ by (ii) at $p = \tfrac23$ (order $\tfrac32$, index $0$). With $x \propto k^{-6}$, $d/d\ln k = -6x\,d/dx$, so $N = -6xG'/G$ and $\alpha = -6x\,dN/dx = 36x\,G'/G + 36x^{2}(G''G - G'^{2})/G^{2}$, which is the stated identity. As $x \to 0$ with $G'(0) \ne 0$, $N = -6G'(0)x + O(x^{2})$ and $\alpha = -6N + O(x^{2})$ (for $G'(0) = 0 \ne G''(0)$ the ratio would be $-12$); a red tilt needs $G'(0) > 0$, and then $\alpha > 0$. The sign of $G'(0)$ is not derived here. At $N = -0.0351$ the leading value is $0.2106$.
\end{proof}
